\chapter{2003年粤桂卷}
\section{单选题}
\begin{problem}
在同一坐标系中，表示直线 \(y = ax\) 与 \(y = x + a\) 正确的是（\quad）

\choices
    {\choicebitmap[width=0.15\paperwidth]{img/2003/yue_gui/q1_optA.png}}
    {\choicebitmap[width=0.15\paperwidth]{img/2003/yue_gui/q1_optB.png}}
    {\choicebitmap[width=0.15\paperwidth]{img/2003/yue_gui/q1_optC.png}}
    {\choicebitmap[width=0.15\paperwidth]{img/2003/yue_gui/q1_optD.png}}
\end{problem}

\begin{answer}
C
\end{answer}

\begin{solution}
直线 \(y=ax\) 必过原点，斜率为 \(a\)；直线 \(y=x+a\) 的斜率为 \(1\)，在 \(y\) 轴上的截距为 \(a\). 当 \(a>0\) 时，两条直线的斜率都为正，且第二条直线在 \(y\) 轴上的截距为正；当 \(a<0\) 时，直线 \(y=ax\) 的斜率为负并过原点，而直线 \(y=x+a\) 的斜率为正且在 \(y\) 轴上的截距为负；当 \(a=0\) 时，两条直线分别为 \(y=0\) 和 \(y=x\)，都过原点.

观察四个图形，只有第三个图形呈现一条负斜率直线过原点、另一条斜率为 \(1\) 且截距为负数的情形，对应 \(a<0\). 故选第三项.
\end{solution}

\begin{problem}
已知 \(x \in \left(-\frac{\pi}{2}, 0\right)\)，\(\cos x = \frac{4}{5}\)，则 \(\tan 2x =\)（\quad）

\choices
    {\(\frac{7}{24}\)}
    {\(-\frac{7}{24}\)}
    {\(\frac{24}{7}\)}
    {\(-\frac{24}{7}\)}
\end{problem}

\begin{answer}
D
\end{answer}

\begin{solution}
因为 \(x\in\left(-\frac{\pi}{2},0\right)\)，所以 \(\sin x<0\). 由 \(\sin ^2x+\cos ^2x=1\)，得
\(\sin x=-\sqrt{1-\frac{16}{25}}=-\frac35\)，\(\tan x=\frac{\sin x}{\cos x}=-\frac34\).
于是
\[
\tan 2x=\frac{2\tan x}{1-\tan ^2x}=\frac{-\frac32}{1-\frac9{16}}=-\frac{24}{7}
\]
故选第四项.
\end{solution}

\begin{problem}
圆锥曲线 \(\rho = \frac{8\sin\theta}{\cos^2\theta}\) 的准线方程是（\quad）

\choices
    {\(\rho \cos \theta = -2\)}
    {\(\rho \cos \theta = 2\)}
    {\(\rho \sin \theta = -2\)}
    {\(\rho \sin \theta = 2\)}
\end{problem}

\begin{answer}
C
\end{answer}

\begin{solution}
在极坐标与直角坐标的关系 \(x=\rho\cos\theta\)，\(y=\rho\sin\theta\) 下，原方程化为
\[
\rho\cos ^2\theta=8\sin\theta
\]
两边同乘 \(\rho\)，并利用 \(\rho^2\cos ^2\theta=x^2\)、\(\rho\sin\theta=y\)，得
\[
x^2=8y
\]
即 \(x^2=4\cdot2\,y\). 这是焦点为 \((0,2)\) 的抛物线，其准线为 \(y=-2\). 写成极坐标方程即
\[
\rho\sin\theta=-2
\]
故选第三项.
\end{solution}

\begin{problem}
等差数列 \(\{a_{n}\}\) 中，已知 \(a_1 = \frac{1}{3}\)，\(a_2 + a_5 = 4\)，\(a_n = 33\)，则 \(n\) 为（\quad）

\choices
    {\(48\)}
    {\(49\)}
    {\(50\)}
    {\(51\)}
\end{problem}

\begin{answer}
C
\end{answer}

\begin{solution}
设等差数列的公差为 \(d\)，则
\[
a_2+a_5=\left(\frac13+d\right)+\left(\frac13+4d\right)=\frac23+5d=4
\]
所以 \(d=\frac23\). 因而
\[
a_n=a_1+(n-1)d=\frac13+\frac23(n-1)=\frac{2n-1}{3}
\]
由 \(a_n=33\)，得 \(2n-1=99\)，所以 \(n=50\). 故选第三项.
\end{solution}

\begin{problem}
双曲线虚轴的一个端点为 \(M\)，两个焦点为 \(F_{1}\)，\(F_{2}\)，\(\angle F_{1}MF_{2} = 120^{\circ}\)，则双曲线的离心率为（\quad）

\choices
    {\(\sqrt{3}\)}
    {\(\frac{\sqrt{6}}{2}\)}
    {\(\frac{\sqrt{6}}{3}\)}
    {\(\frac{\sqrt{3}}{3}\)}
\end{problem}

\begin{answer}
B
\end{answer}

\begin{solution}
设双曲线的实半轴、虚半轴和半焦距分别为 \(a\)，\(b\)，\(c\)，则 \(c^2=a^2+b^2\). 取坐标系使 \(M=(0,b)\)，\(F_1=(-c,0)\)，\(F_2=(c,0)\)，于是
\(MF_1^2=MF_2^2=b^2+c^2\)，\(F_1F_2=2c\).
在三角形 \(F_1MF_2\) 中，由余弦定理
\[
\cos120^{\circ}=\frac{2(b^2+c^2)-4c^2}{2(b^2+c^2)}=\frac{b^2-c^2}{b^2+c^2}=-\frac12
\]
整理得 \(c^2=3b^2\)，从而 \(a^2=c^2-b^2=2b^2\). 因此离心率
\[
e=\frac ca=\frac{\sqrt3b}{\sqrt2b}=\frac{\sqrt6}{2}
\]
故选第二项.
\end{solution}

\begin{problem}
设函数 \( f(x) = \begin{cases}
2^{-x} - 1, & x \leqslant 0,\\
x^{\frac{1}{2}}, & x > 0,
\end{cases} \) 若 \( f(x_0) > 1 \)，则 \( x_0 \) 的取值范围是（\quad）

\choices
    {\((-1, 1)\)}
    {\((-1, +\infty)\)}
    {\((-\infty, -2) \cup (0, +\infty)\)}
    {\((-\infty, -1) \cup (1, +\infty)\)}
\end{problem}

\begin{answer}
D
\end{answer}

\begin{solution}
当 \(x_0\leqslant0\) 时，\(f(x_0)=2^{-x_0}-1\). 由 \(f(x_0)>1\)，得 \(2^{-x_0}>2\)，所以 \(-x_0>1\)，即 \(x_0<-1\). 这与 \(x_0\leqslant0\) 相容.

当 \(x_0>0\) 时，\(f(x_0)=\sqrt{x_0}\). 由 \(\sqrt{x_0}>1\)，得 \(x_0>1\).

综上，\(x_0\) 的取值范围为 \((-\infty,-1)\cup(1,+\infty)\). 故选第四项.
\end{solution}

\begin{problem}
函数 \(y = 2\sin x (\sin x + \cos x)\) 的最大值为（\quad）

\choices
    {\(1 + \sqrt{2}\)}
    {\(\sqrt{2} - 1\)}
    {\(\sqrt{2}\)}
    {\(2\)}
\end{problem}

\begin{answer}
A
\end{answer}

\begin{solution}
利用二倍角公式，得
\[
y=2\sin ^2x+2\sin x\cos x=1-\cos2x+\sin2x
  =1+\sqrt2\sin\left(2x-\frac\pi4\right)
\]
因为正弦函数的最大值为 \(1\)，所以 \(y\) 的最大值为 \(1+\sqrt2\). 故选第一项.
\end{solution}

\begin{problem}
已知圆 \(C: (x - a)^2 + (y - 2)^2 = 4\)（\(a > 0\)）及直线 \(l: x - y + 3 = 0\)，当直线 \(l\) 被 \(C\) 截得的弦长为 \(2\sqrt{3}\) 时，\(a\) 的值等于（\quad）

\choices
    {\(\sqrt{2}\)}
    {\(2 - \sqrt{2}\)}
    {\(\sqrt{2} - 1\)}
    {\(\sqrt{2} + 1\)}
\end{problem}

\begin{answer}
C
\end{answer}

\begin{solution}
圆 \(C\) 的半径为 \(R=2\)，圆心为 \(M(a,2)\). 设圆心到直线 \(l\) 的距离为 \(d\)，则弦长公式给出
\[
2\sqrt{R^2-d^2}=2\sqrt3
\]
所以 \(4-d^2=3\)，即 \(d=1\). 另一方面，
\[
d=\frac{|a-2+3|}{\sqrt{1^2+(-1)^2}}=\frac{a+1}{\sqrt2}
\]
其中利用了 \(a>0\). 因而 \(a+1=\sqrt2\)，得 \(a=\sqrt2-1\). 故选第三项.
\end{solution}

\begin{problem}
已知圆锥的底面半径为 \(R\)，高为 \(3R\)，在它的所有内接圆柱中，全面积的最大值是（\quad）

\choices
    {\(2\pi R^2\)}
    {\(\frac{9}{4}\pi R^2\)}
    {\(\frac{8}{3}\pi R^2\)}
    {\(\frac{3}{2}\pi R^2\)}
\end{problem}

\begin{answer}
B
\end{answer}

\begin{solution}
设内接圆柱的底面半径为 \(r\)，高为 \(h\). 由相似三角形，
\(\frac rR=\frac{3R-h}{3R}\)，\(h=3(R-r)\)，\(0\leqslant r\leqslant R\).
圆柱的全面积为
\[
S=2\pi r^2+2\pi rh=2\pi r\bigl(r+3R-3r\bigr)=6\pi Rr-4\pi r^2
\]
这是关于 \(r\) 的开口向下二次函数，其顶点在 \(r=\frac{3R}{4}\)，且该值属于 \([0,R]\). 因此
\[
S_{\max}=6\pi R\cdot\frac{3R}{4}-4\pi\left(\frac{3R}{4}\right)^2=\frac94\pi R^2
\]
故选第二项.
\end{solution}

\begin{problem}
函数 \(f(x) = \sin x\)，\(x \in \left[\frac{\pi}{2}, \frac{3\pi}{2}\right]\) 的反函数 \(f^{-1}(x) =\)（\quad）

\choices
    {\(-\arcsin x,\ x \in [-1,1]\)}
    {\(-\pi - \arcsin x\)，\(\ x \in [-1,1]\)}
    {\(\pi + \arcsin x\)，\(\ x \in [-1,1]\)}
    {\(-\pi - \arcsin x\)，\(\ x \in [-1,1]\)}
\end{problem}

\begin{answer}
原题无正确选项
\end{answer}

\begin{solution}
在区间 \(\left[\frac\pi2,\frac{3\pi}2\right]\) 上，\(\sin x\) 严格递减，值域为 \([-1,1]\)，所以反函数的定义域为 \([-1,1]\)，值域为 \(\left[\frac\pi2,\frac{3\pi}2\right]\). 令 \(y=f^{-1}(x)\)，则 \(y\in\left[\frac\pi2,\frac{3\pi}2\right]\) 且 \(\sin y=x\). 由 \(\arcsin x\in\left[-\frac\pi2,\frac\pi2\right]\)，角 \(\pi-\arcsin x\) 属于 \(\left[\frac\pi2,\frac{3\pi}2\right]\)，并且
\[
y=\pi-\arcsin x
\]
因此按函数定义，反函数应为 \(f^{-1}(x)=\pi-\arcsin x\)，其中 \(x\in[-1,1]\). 直接代入端点也可核验：当 \(x=1\) 时，\(f^{-1}(1)=\frac\pi2\)；当 \(x=-1\) 时，\(f^{-1}(-1)=\frac{3\pi}2\).

原始试卷的第二、第四个选项均印为 \(-\pi-\arcsin x\)，没有印出上述正确表达式，因此原题无正确选项；原卷参考答案记为第四项，但该答案与题面及反函数定义不一致.
\end{solution}

\begin{problem}
已知长方形的四个顶点 \(A(0,0)\)，\(B(2,0)\)，\(C(2,1)\) 和 \(D(0,1)\)，一质点从 \(AB\) 的中点 \(P_{0}\) 沿与 \(AB\) 的夹角 \(\theta\) 的方向射到 \(BC\) 上的点 \(P_{1}\) 后，依次反射到 \(CD\)，\(DA\) 和 \(AB\) 上的点 \(P_{2}\)，\(P_{3}\) 和 \(P_{4}\)（入射角等于反射角）. 设 \(P_{4}\) 的坐标为 \((x_{4},0)\)，若 \(1 < x_{4} < 2\)，则 \(\tan \theta\) 的取值范围是（\quad）

\choices
    {\(\left(\frac{1}{3}, 1\right)\)}
    {\(\left(\frac{1}{3}, \frac{2}{3}\right)\)}
    {\(\left(\frac{2}{5}, \frac{1}{2}\right)\)}
    {\(\left(\frac{2}{5}, \frac{3}{5}\right)\)}
\end{problem}

\begin{answer}
C
\end{answer}

\begin{solution}
设 \(m=\tan\theta\). 将质点每次反射后的矩形关于反射边展开，则折线路径变为一条直线. 在展开图中，取 \(P_0=(1,0)\)，直线方程为
\[
y=m(x-1)
\]
依次经过右边、上边、左边时，展开图中相应的边界依次为 \(x=2\)、\(y=1\)、\(x=4\)，第四次到达的底边对应于 \(y=2\). 因此直线先后与前三条边相交的坐标分别满足
\(m<1\)，\(2<1+\frac1m<4\)，\(1<3m<2\).
第四次交点的横坐标为
\[
x'=1+\frac2m
\]
要使第四次交点位于左边界反射后的矩形内，还需有 \(4<x'<6\). 反射回原矩形后，\(P_4\) 的横坐标为
\[
x_4=x'-4=\frac2m-3
\]
题设 \(1<x_4<2\) 等价于
\[
1<\frac2m-3<2
\]
从而
\[
\frac25<m<\frac12
\]
这个范围同时满足前述各边界相交条件，所以直线确实依次穿过右边、上边、左边和底边，且不经过顶点. 因此
\[
\tan\theta\in\left(\frac25,\frac12\right)
\]
故选第三项.
\end{solution}

\begin{problem}
一个四面体的所有棱长都为 \(\sqrt{2}\)，四个顶点在同一球面上，则此球的表面积为（\quad）

\choices
    {\(3\pi\)}
    {\(4\pi\)}
    {\(3\sqrt{3}\pi\)}
    {\(6\pi\)}
\end{problem}

\begin{answer}
A
\end{answer}

\begin{solution}
这是棱长为 \(a=\sqrt2\) 的正四面体. 设顶点到底面的高为 \(h\)，底面中心到任一底面顶点的距离为 \(\frac{a}{\sqrt3}\)，则
\[
h=\sqrt{a^2-\left(\frac{a}{\sqrt3}\right)^2}=a\sqrt{\frac23}
\]
设球心到对面底面中心的距离为 \(u\)，则球心到顶点的距离为 \(h-u\)，到任一底面顶点的距离满足
\[
(h-u)^2=u^2+\frac{a^2}{3}
\]
代入 \(h^2=\frac23a^2\)，得 \(u=\frac h4\)，所以外接球半径为
\[
R=\frac34a\sqrt{\frac23}=\frac{a\sqrt6}{4}=\frac{\sqrt3}{2}
\]
因此球的表面积为
\[
4\pi R^2=4\pi\cdot\frac34=3\pi
\]
故选第一项.
\end{solution}

\section{填空题}
\begin{problem}
不等式 \(\sqrt{4x - x^{2}} < x\) 的解集是\fillinblank{}.
\end{problem}

\begin{answer}
\((2,4]\)
\end{answer}

\begin{solution}
由根式有意义得 \(4x-x^{2}\geqslant0\)，所以 \(0\leqslant x\leqslant4\). 因为左边非负，而题设要求左边小于 \(x\)，故必须有 \(x>0\). 在此条件下两边均非负，可以平方，得 \(4x-x^{2}<x^{2}\)，即 \(2x(x-2)>0\)，所以 \(x>2\). 与定义域合并，得原不等式的解集为 \((2,4]\).
\end{solution}

\begin{problem}
\(\left(x^{2} - \frac{1}{2x}\right)^{9}\) 的展开式中 \(x^{9}\) 系数是\fillinblank{}.
\end{problem}

\begin{answer}
\(-\frac{21}{2}\)
\end{answer}

\begin{solution}
展开式通项中取 \(\left(-\frac{1}{2x}\right)^{k}\) 共 \(k\) 次时，该项为
\(\mathrm C_{9}^{k}\left(x^{2}\right)^{9-k}\left(-\frac{1}{2x}\right)^{k}=\mathrm C_{9}^{k}\frac{(-1)^{k}}{2^{k}}x^{18-3k}\)，\(k=0\)，\(1\)，\(\ldots\)，\(9\).
要得到 \(x^{9}\)，必须满足 \(18-3k=9\)，解得 \(k=3\). 因此所求系数为
\[
\mathrm C_{9}^{3}\frac{(-1)^{3}}{2^{3}}=-\frac{84}{8}=-\frac{21}{2}
\]
\end{solution}

\begin{problem}
在平面几何里，有勾股定理：“设 \(\triangle ABC\) 的两边 \(AB\)，\(AC\) 互相垂直，则 \(AB^2 + AC^2 = BC^2\).”拓展到空间，类比平面几何的勾股定理，研究三棱锥的侧面面积与底面面积间的关系，可以得出的正确结论是：“设三棱锥 \(A-BCD\) 的三个侧面 \(ABC\)，\(ACD\)，\(ADB\) 两两互相垂直，则\fillinblank{}.”

\end{problem}

\begin{answer}
\(S_{\triangle ABC}^{2}+S_{\triangle ACD}^{2}+S_{\triangle ADB}^{2}=S_{\triangle BCD}^{2}\)
\end{answer}

\begin{solution}
因为平面 \(ABC\)，\(ACD\)，\(ADB\) 两两互相垂直，且它们的交棱分别为 \(AC\)，\(AD\)，\(AB\)，可取直角坐标系使 \(A=(0,0,0)\)，\(B=(u,0,0)\)，\(C=(0,v,0)\)，\(D=(0,0,w)\)，其中 \(u=AB>0\)，\(v=AC>0\)，\(w=AD>0\).

事实上，先取 \(AC\) 为 \(x\) 轴，使平面 \(ABC\)，\(ACD\) 分别为 \(xy\) 平面、\(xz\) 平面，设 \(B=(b_1,b_2,0)\)，\(D=(d_1,0,d_3)\). 平面 \(ABD\) 的一个法向量为 \((b_2d_3,-b_1d_3,-b_2d_1)\)，而平面 \(ABC\)，\(ACD\) 的法向量分别可取 \((0,0,1)\)，\((0,1,0)\). 由平面 \(ABD\) 分别垂直于 \(ABC\)、\(ACD\)，以及 \(b_2\ne0\)，\(d_3\ne0\)，得 \(d_1=0\)，\(b_1=0\)，所以 \(AB\perp AC\)，\(AD\perp AC\)，且 \(AB\perp AD\)，上述坐标表示成立.

于是 \(S_{\triangle ABC}=\frac{uv}{2}\)，\(S_{\triangle ACD}=\frac{vw}{2}\)，\(S_{\triangle ADB}=\frac{wu}{2}\).

又 \(\overrightarrow{BC}=(-u,v,0)\)，\(\overrightarrow{BD}=(-u,0,w)\).

从而

\[
\overrightarrow{BC}\times\overrightarrow{BD}=(vw,uw,uv)
\]

因此

\[
S_{\triangle BCD}^{2}=\frac14\left(u^{2}v^{2}+v^{2}w^{2}+w^{2}u^{2}\right)=S_{\triangle ABC}^{2}+S_{\triangle ACD}^{2}+S_{\triangle ADB}^{2}
\]

所以

\[
S_{\triangle ABC}^{2}+S_{\triangle ACD}^{2}+S_{\triangle ADB}^{2}=S_{\triangle BCD}^{2}
\]
\end{solution}

\begin{problem}
如图，一个地区分为 \(5\) 个行政区域，现给地图着色，要求相邻地区不得使用同一颜色，现有 \(4\) 种颜色可供选择，则不同的着色方法共有\fillinblank{}种.（以数字作答）

\bitmapfigure[max width=0.18\linewidth]{img/2003/yue_gui/q16_fig1.png}

\end{problem}

\begin{answer}
72
\end{answer}

\begin{solution}
区域 \(1\) 与区域 \(2\)，\(3\)，\(4\)，\(5\) 都相邻. 先给区域 \(1\) 选颜色，有 \(4\) 种选法；其余四个区域只能使用剩下的 \(3\) 种颜色.

从图中可知，区域 \(2\)，\(5\)，\(4\)，\(3\) 依次相邻并首尾相邻，构成一个四边形环. 固定区域 \(2\) 的颜色后，区域 \(5\) 有 \(2\) 种选法. 对每一种区域 \(2\)、\(5\) 的不同颜色组合，设剩余的第三种颜色为 \(c\). 区域 \(3\) 不能与区域 \(2\) 同色，因此有两种可能：

当区域 \(3\) 与区域 \(5\) 同色时，区域 \(4\) 只需避开这一种颜色，有 \(2\) 种选法；当区域 \(3\) 使用第三种颜色 \(c\) 时，区域 \(4\) 必须使用区域 \(2\) 的颜色，有 \(1\) 种选法. 所以固定区域 \(2\) 的颜色后有 \(2(2+1)=6\) 种着色方法，区域 \(2\) 有 \(3\) 种颜色可选，四个外部区域共有 \(3\times6=18\) 种着色方法.

因此不同的着色方法共有 \(4\times18=72\) 种.
\end{solution}

\section{解答题}
\begin{problem}
已知正四棱柱 \(ABCD-A_{1}B_{1}C_{1}D_{1}\)，\(AB = 1\)，\(AA_{1} = 2\)，\(E\) 为 \(CC_{1}\) 中点，\(F\) 为 \(BD_{1}\) 中点.

\begin{enumerate}
\item 证明：\(EF\) 为 \(BD_{1}\) 与 \(CC_{1}\) 的公垂线；
\item 求点 \(D_{1}\) 到面 \(BDE\) 的距离.
\end{enumerate}

\bitmapfigure[max width=0.30\linewidth]{img/2003/yue_gui/q17_fig1.png}
\end{problem}

\begin{solution}
\begin{enumerate}
\item 以 \(A\) 为原点，分别以 \(\overrightarrow{AB}\)，\(\overrightarrow{AD}\)，\(\overrightarrow{AA_{1}}\) 的方向为 \(x\)，\(y\)，\(z\) 轴正方向，建立空间直角坐标系，则
\(B(1,0,0)\)，\(D(0,1,0)\)，\(C(1,1,0)\)，\(D_{1}(0,1,2)\)，\(C_{1}(1,1,2)\).
由中点坐标公式，得
\(E(1,1,1)\)，\(F\left(\frac12,\frac12,1\right)\).
于是
\(\overrightarrow{EF}=\left(-\frac12,-\frac12,0\right)\)，\(\overrightarrow{BD_{1}}=(-1,1,2)\)，\(\overrightarrow{CC_{1}}=(0,0,2)\).
从而
\(\overrightarrow{EF}\cdot\overrightarrow{BD_{1}}=0\)，\(\overrightarrow{EF}\cdot\overrightarrow{CC_{1}}=0\).
又 \(F\in BD_{1}\)，\(E\in CC_{1}\)，所以 \(EF\) 同时垂直于 \(BD_{1}\) 和 \(CC_{1}\)，即 \(EF\) 是这两条异面直线的公垂线.
\item 由
\(\overrightarrow{BD}=(-1,1,0)\)，\(\overrightarrow{BE}=(0,1,1)\)
得平面 \(BDE\) 的一个法向量为
\[
\overrightarrow{BD}\times\overrightarrow{BE}=(1,1,-1)
\]
因此平面 \(BDE\) 的方程为
\[
(x-1)+y-z=0
\]
即 \(x+y-z-1=0\). 点 \(D_{1}(0,1,2)\) 到该平面的距离为
\[
\frac{|0+1-2-1|}{\sqrt{1^{2}+1^{2}+(-1)^{2}}}
=\frac{2}{\sqrt3}=\frac{2\sqrt3}{3}
\]
\end{enumerate}
\end{solution}

\begin{problem}
已知复数 \(z\) 的辐角为 \(60^{\circ}\)，且 \(|z - 1|\) 是 \(|z|\) 和 \(|z - 2|\) 的等比中项，求 \(|z|\).
\end{problem}

\begin{solution}
设 \(r=|z|\). 因为 \(z\) 的辐角为 \(60^{\circ}\)，所以 \(r>0\)，且
\[
z=r\left(\cos60^{\circ}+\i\sin60^{\circ}\right)
=\frac r2+\frac{\sqrt3}{2}r\i
\]
由“\(|z-1|\) 是 \(|z|\) 和 \(|z-2|\) 的等比中项”，得
\[
|z-1|^{2}=|z|\,|z-2|
\]
分别计算两边所需的模：
\[
|z-1|^{2}=\left(\frac r2-1\right)^{2}+\left(\frac{\sqrt3}{2}r\right)^{2}
=r^{2}-r+1
\]
\[
|z-2|^{2}=\left(\frac r2-2\right)^{2}+\left(\frac{\sqrt3}{2}r\right)^{2}
=r^{2}-2r+4
\]
故
\[
r^{2}-r+1=r\sqrt{r^{2}-2r+4}
\]
两边均为非负数，平方得
\[
(r^{2}-r+1)^{2}=r^{2}(r^{2}-2r+4)
\]
展开并约去 \(r^{4}-2r^{3}\)，得到
\[
r^{2}+2r-1=0
\]
所以
\[
r=-1\pm\sqrt2
\]
由于 \(r>0\)，只能取 \(r=\sqrt2-1\). 代回原等式时两边均为正，故该值成立，从而
\[
|z|=\sqrt2-1
\]
\end{solution}

\begin{problem}
已知 \(c > 0\)，设 \(P\)：函数 \(y = c^{x}\) 在 \(\R\) 上单调递减；\(Q\)：不等式 \(x + |x - 2c| > 1\) 的解集为 \(\R\). 如果 \(P\) 和 \(Q\) 有且仅有一个正确，求 \(c\) 的取值范围.
\end{problem}

\begin{solution}
因为指数函数 \(y=c^{x}\) 在 \(\R\) 上单调递减当且仅当 \(0<c<1\)，所以命题 \(P\) 正确的充要条件是 \(0<c<1\). 对命题 \(Q\)，有
\[
x+|x-2c|=
\begin{cases}
2c, & x\leqslant2c,\\
2x-2c, & x>2c.
\end{cases}
\]
因此该式的最小值为 \(2c\). 不等式对一切实数 \(x\) 都成立的充要条件是 \(2c>1\)，即命题 \(Q\) 正确的充要条件是 \(c>\frac12\).

当 \(0<c<1\) 时，\(P\) 正确；要使 \(P\) 和 \(Q\) 恰有一个正确，必须 \(Q\) 错误，即 \(0<c\leqslant\frac12\). 当 \(c\geqslant1\) 时，\(P\) 错误而 \(Q\) 正确. 因此 \(c\) 的取值范围为
\[
\left(0,\frac12\right]\cup[1,+\infty)
\]
\end{solution}

\begin{problem}
在某海滨城市附近海面有一台风，据监测，当前台风中心位于城市 \(O\)（如图）的东偏南 \(\theta\)（\(\theta = \arccos \frac{\sqrt{2}}{10}\)）方向 \(300 \mathrm{~km}\) 的海面 \(P\) 处，并以 \(20 \mathrm{~km}/\mathrm{h}\) 的速度向西偏北 \(45^{\circ}\) 方向移动，台风侵袭的范围为圆形区域. 当前半径为 \(60 \mathrm{~km}\)，并以 \(10 \mathrm{~km}/\mathrm{h}\) 的速度不断增大，问几小时后该城市开始受到台风的侵袭？

\bitmapfigure[width=0.80\linewidth]{img/2003/yue_gui/q20_fig1.png}
\end{problem}

\begin{solution}
以城市 \(O\) 为原点，向东、向北分别为 \(x\)，\(y\) 轴正方向，并设 \(t\) 为经过的时间，单位为 \(\mathrm h\). 因为台风中心初始位于 \(O\) 的东偏南 \(\theta\) 方向，且
\(\cos\theta=\frac{\sqrt2}{10}\)，\(\sin\theta=\frac{7\sqrt2}{10}\)
所以初始位置为
\[
P_{0}=(300\cos\theta,-300\sin\theta)=(30\sqrt2,-210\sqrt2)
\]
台风以 \(20\mathrm{~km}/\mathrm h\) 的速度向西偏北 \(45^{\circ}\) 方向移动，故 \(t\) 小时后的位移为
\[
20t\left(-\cos45^{\circ},\sin45^{\circ}\right)=(-10\sqrt2t,10\sqrt2t)
\]
于是台风中心的位置为
\[
P(t)=\left(30\sqrt2-10\sqrt2t,-210\sqrt2+10\sqrt2t\right)
\]
城市开始受到侵袭时，台风中心到城市的距离等于台风半径. 由上式，中心到 \(O\) 的距离平方为
\[
d(t)^{2}=200\left[(t-3)^{2}+(t-21)^{2}\right]
\]
而台风半径为
\[
r(t)=60+10t
\]
因此接触时刻满足
\[
200\left[(t-3)^{2}+(t-21)^{2}\right]=(60+10t)^{2}
\]
即
\[
t^{2}-36t+288=0
\]
解得 \(t=12\) 或 \(t=24\). 开始受到侵袭对应第一次接触，所以取较小的非负根，得 \(t=12\). 即该城市在 \(12 \mathrm{~h}\) 后开始受到台风侵袭.
\end{solution}

\begin{problem}
已知常数 \(a > 0\)，在矩形 \(ABCD\) 中，\(AB = 4\)，\(BC = 4a\)，\(O\) 为 \(AB\) 的中点. 点 \(E\)，\(F\)，\(G\) 分别在 \(BC\)，\(CD\)，\(DA\) 上移动. 且 \(\frac{BE}{BC} = \frac{CF}{CD} = \frac{DG}{DA}\)，\(P\) 为 \(GE\) 与 \(OF\) 的交点（如图），问是否存在两个定点，使 \(P\) 到这两点的距离的和为定值？若存在，求出这两点的坐标及此定值；若不存在，请说明理由.

\bitmapfigure[max width=0.35\linewidth]{img/2003/yue_gui/q21_fig1.png}
\end{problem}

\begin{solution}
建立如图所示的直角坐标系，取 \(O\) 为原点，\(x\) 轴沿 \(AB\) 的方向，\(y\) 轴沿 \(BC\) 的方向，则
\(A(-2,0)\)，\(B(2,0)\)，\(C(2,4a)\)，\(D(-2,4a)\).
设公共比值为 \(\lambda\)，则 \(0\leqslant\lambda\leqslant1\)，从而
\(E=(2,4a\lambda)\)，\(F=(2-4\lambda,4a)\)，\(G=(-2,4a(1-\lambda))\).
令 \(s=1-2\lambda\)，则 \(-1\leqslant s\leqslant1\)，且
\[
F=(2s,4a)
\]
直线 \(OF\) 上的点可表示为 \(P=(2s\mu,4a\mu)\). 另一方面，由 \(G\)，\(E\) 两点的坐标，直线 \(GE\) 的方程为
\[
y=2a-asx
\]
将 \(P\) 的坐标代入直线 \(GE\)，得
\[
4a\mu=2a-2as^{2}\mu
\]
所以
\[
\mu=\frac{1}{s^{2}+2}
\]
因此点 \(P\) 的坐标为
\[
P\left(\frac{2s}{s^{2}+2},\frac{4a}{s^{2}+2}\right)
\]
消去参数 \(s\)，可得
\[
x^{2}+\frac{(y-a)^{2}}{2a^{2}}=\frac12
\]
这是以 \((0,a)\) 为中心的椭圆. 当 \(a>\frac{1}{\sqrt2}\) 时，其长轴在 \(y\) 轴上，两个焦点为
\(F_{1}\left(0,a-\sqrt{a^{2}-\frac12}\right)\)，\(F_{2}\left(0,a+\sqrt{a^{2}-\frac12}\right)\)
到这两个焦点的距离和为 \(2a\). 当 \(0<a<\frac{1}{\sqrt2}\) 时，其长轴在 \(x\) 轴上，两个焦点为
\(F_{1}\left(-\sqrt{\frac12-a^{2}},a\right)\)，\(F_{2}\left(\sqrt{\frac12-a^{2}},a\right)\)
到这两个焦点的距离和为 \(\sqrt2\). 当 \(a=\frac{1}{\sqrt2}\) 时，轨迹方程为
\[
x^{2}+\left(y-\frac{1}{\sqrt2}\right)^{2}=\frac12
\]
是圆弧. 若存在两个不同定点，使 \(P\) 到这两点的距离和为定值，则这段圆弧应属于以这两点为焦点的椭圆；两个焦点不重合的椭圆不能退化为圆，故不存在这样的两个定点. 因此，当
\[
a\in\left(0,\frac{1}{\sqrt2}\right)\cup\left(\frac{1}{\sqrt2},+\infty\right)
\]
时存在这样的两个定点；当 \(a=\frac{1}{\sqrt2}\) 时不存在.
\end{solution}

\begin{problem}
设 \(a_0\) 为常数，且 \(a_n = 3^{n-1} - 2a_{n-1}\)（\(n \in \N^*\)）.

\begin{enumerate}
\item 证明对任意 \(n \geqslant 1\)，\(a_n = \frac{1}{5}\left(3^n + (-1)^{n-1} \cdot 2^n\right) + (-1)^n \cdot 2^n a_0\)；
\item 假设对任意 \(n \geqslant 1\) 有 \(a_n > a_{n-1}\)，求 \(a_0\) 的取值范围.
\end{enumerate}
\end{problem}

\begin{solution}
\begin{enumerate}
\item 当 \(n=1\) 时，由递推关系得 \(a_{1}=1-2a_{0}\)，而题中公式右端为
\[
\frac15\left(3+(-1)^{0}\cdot2\right)+(-1)^{1}\cdot2a_{0}=1-2a_{0}
\]
所以公式在 \(n=1\) 时成立. 假设当 \(n=k\) 时成立，即
\[
a_{k}=\frac15\left(3^{k}+(-1)^{k-1}\cdot2^{k}\right)+(-1)^{k}\cdot2^{k}a_{0}
\]
则由递推关系
\[
a_{k+1}=3^{k}-2a_{k}
=3^{k}-\frac25\left(3^{k}+(-1)^{k-1}\cdot2^{k}\right)-(-1)^{k}\cdot2^{k+1}a_{0}
\]
由于 \(-(-1)^{k-1}=(-1)^k\)，且 \(3^k-\frac25\cdot3^k=\frac15\cdot3^{k+1}\)，上式化为
\[
a_{k+1}=\frac15\left(3^{k+1}+(-1)^{k}\cdot2^{k+1}\right)+(-1)^{k+1}\cdot2^{k+1}a_{0}
\]
这正是 \(n=k+1\) 时的结论，故由数学归纳法，对任意 \(n\geqslant1\)，均有
\[
a_{n}=\frac15\left(3^{n}+(-1)^{n-1}\cdot2^{n}\right)+(-1)^{n}\cdot2^{n}a_{0}
\]
\item 由上式，得
\[
a_{n}-a_{n-1}=\frac{2\cdot3^{n-1}}5+3\cdot2^{n-1}(-1)^{n-1}\left(\frac15-a_{0}\right)
\]
当 \(n=2k+1\)（\(k\geqslant0\)）时，
\[
a_{2k+1}-a_{2k}
=\frac{2\cdot3^{2k}}5+3\cdot2^{2k}\left(\frac15-a_{0}\right)>0
\]
等价于
\[
a_{0}<\frac15+\frac{2}{15}\left(\frac94\right)^{k}
\]
因为 \(\left(\frac94\right)^{k}\geqslant1\)，所有奇数项不等式成立的充要条件为 \(a_{0}<\frac13\). 当 \(n=2k\)（\(k\geqslant1\)）时，
\[
a_{2k}-a_{2k-1}
=\frac{2\cdot3^{2k-1}}5-3\cdot2^{2k-1}\left(\frac15-a_{0}\right)>0
\]
等价于
\[
a_{0}>\frac15-\frac{2}{15}\left(\frac32\right)^{2k-1}
\]
右端在 \(k=1\) 时最大且等于 \(0\)，所以所有偶数项不等式成立的充要条件为 \(a_{0}>0\). 综合两类条件，得
\[
0<a_{0}<\frac13
\]
\end{enumerate}
\end{solution}
